% ============================================================
% CHAPTER — Conclusion and Future Work (NEW)
%
% This chapter did not exist before; "Revisiting the Research
% Questions" and the MRQ synthesis are moved here verbatim in
% substance (tightened into a table + short prose) from the old
% closing section of Chapter_Discussion_REVISED.tex, since
% Discussion is no longer the thesis's final chapter. The
% three-item contribution list previously embedded in
% §discussion_roles of Discussion is restated here in full as the
% canonical C1--C3 statement. Future Work is new content, derived
% directly from the Limitations section of Chapter~\ref{chap:discussion}.
%
% Label sec:eq_revisited is defined here (moved from Discussion);
% Chapter5_ExperimentalSetup_REVISED.tex references it twice and
% does not care which chapter file defines it.
% ============================================================

\chapter{Conclusion and Future Work}
\label{chap:conclusion}

This chapter closes the thesis. Section~\ref{sec:conclusion_summary} briefly recaps what was built and found. Section~\ref{sec:eq_revisited} answers each sub-research question stated in Section~\ref{sec:research_questions} and synthesises the answer to the main research question, completing the traceability chain introduced in Table~\ref{tab:rq_traceability}. Section~\ref{sec:contributions} restates the thesis's contributions in full. Section~\ref{sec:future_work} turns the limitations identified in Section~\ref{sec:limitations} into concrete directions for future work.

% ------------------------------------------------------------
\section{Summary}
\label{sec:conclusion_summary}

This thesis set out to answer whether reinforcement learning can optimise declarative business-process execution while guaranteeing compliance with \ac{DCR} constraints. The proposed framework models \ac{DCR} execution as a shielded sequential decision problem, scalarises competing cost and duration objectives into a single reward signal, and extends the action space to joint event--role decisions for resource allocation. On the Expense Report benchmark, the trained \ac{PPO} agent recovers the exact two-point Pareto front identified by a CSP+COP ground-truth solver (Chapter~\ref{chap:results}), while remaining fully compliant from the first training episode. On the larger Loan Application benchmark, the same approach produces a well-populated 26-point Pareto front and learns event-level role-specialisation patterns that respond systematically to the objective weights. An ablation study (Chapter~\ref{chap:internal_validation}) attributes this performance to the combination of shielding and learning specifically, rather than to either component alone, and the same qualitative findings replicate on six independently mined \ac{DCR} graphs never used to design the method (Chapter~\ref{chap:external_validation}), which also expose the scale, approximately 42 event--role pairs, beyond which the unmasked policy's own search becomes the limiting factor rather than process compliance.

% ------------------------------------------------------------
\section{Revisiting the Research Questions}
\label{sec:eq_revisited}

\paragraph{SRQ1 --- How can \ac{DCR} graph semantics be represented within a \ac{RL} environment?} The answer is constructive. \ac{DCR} execution is modelled as a sequential decision-making problem in which states encode the marking of the graph, actions correspond to event (or event--role) executions, and the \ac{DCR} engine itself acts as the transition function (Sections~\ref{sec:compliance}, \ref{sec:ppo}). The fidelity of this representation is verified empirically through EQ2: the 100\,\% acceptance rate observed from the first episode (Sections~\ref{sec:expense_compliance}, \ref{sec:loanapp_convergence}) confirms that the environment exposes exactly the reachable state space of the underlying model, no more and no less.

\paragraph{SRQ2 --- To what extent can \ac{PPO}-based agents discover Pareto-optimal traces in declarative process models?} On the Expense Report benchmark, the agent recovers the exact two-point Pareto front identified by the CSP+COP ground truth (Section~\ref{sec:pareto_recovery}), answering EQ1 affirmatively. On the larger Loan Application benchmark, where no ground truth exists, the weight sweep produces a 26-point front whose structure is consistent with independent approaches (Section~\ref{sec:loanapp_pareto}). The extent of discovery is bounded by the training budget, the convexity limitation of linear scalarisation (Section~\ref{sec:limitations}), and, at larger scale, by the action-space ceiling established in Chapter~\ref{chap:external_validation}: exact recovery is confirmed up to 42 actions, and the unmasked policy does not discover any accepting trace beyond that range, even though a valid one is known to exist (Section~\ref{sec:external_scalability}). The ablation of Section~\ref{sec:shield_contribution} further shows that this discovery is attributable to the combination of shielding and learning specifically: the random shielded baseline alone recovers a comparable but less stable and non-steerable front, and the unshielded learner recovers no valid front at all.

\paragraph{SRQ3 --- How does reward design influence convergence and compliant agent behaviour?} Reward design proved decisive for convergence. The step penalty competes with the terminal acceptance bonus ($+100$) at every step, and that imbalance, not any explicit cost or duration term, is what structures the trade-off: it makes shorter traces cheaper to execute in reward terms regardless of objective weights, which is why the duration-focused configuration recovers both Expense Report Pareto points while the cost-focused one, receiving no comparable pressure on episode length, does not (Section~\ref{sec:discussion_tradeoff}). The objective weights then steer \emph{which} short trace is preferred, and the balanced configuration acts as an implicit diversity mechanism on top of this, contributing 18 of the 26 Pareto points on the Loan Application. Regarding compliant behaviour, compliance itself is structural rather than learned; what the reward shapes is the agent's reliance on the shield, measured as the illegal-action ratio, which falls below 50\,\% within the first 1\,\% of the training budget (Table~\ref{tab:compliance_milestones}), and the same gap between shielded and unshielded reliance on the reward signal replicates on an external benchmark never used to design the reward (Section~\ref{sec:external_compliance_replication}).

\paragraph{SRQ4 --- Can the framework learn resource-allocation strategies through role-conditioned action spaces?} Yes, within the scale tested. The role-conditioned extension learns event-level specialisation patterns that respond systematically to the objective weights (Section~\ref{sec:loanapp_roles}), exposing three behavioural categories, high-sensitivity, moderate-response, and role-indifferent events (Section~\ref{sec:discussion_roles}). This granularity is unavailable to pool-level approaches and constitutes the empirical basis of contribution~C2. Whether this extends to allocation problems with substantially larger action spaces is constrained by the same ceiling identified under SRQ2.

\begin{table}[H]
\centering
\caption{Summary of how each thesis component answers the main research question (MRQ).}
\label{tab:mrq_answers}
\renewcommand{\arraystretch}{1.3}
\begin{tabularx}{\linewidth}{p{3.4cm} X p{2.6cm}}
\toprule
\textbf{Component} & \textbf{How it answers the MRQ} & \textbf{Evidence} \\
\midrule

Shielded \ac{DCR} environment &
Compliance is guaranteed by construction, not learned: invalid actions are blocked before execution, so every generated trace satisfies the \ac{DCR} constraints throughout training, not only at convergence. &
Ch.~\ref{chap:results}, \ref{chap:discussion} \\
\addlinespace

Scalarised multi-objective reward &
A single weight sweep over $(\alpha,\beta)$ produces a full family of policies covering the cost--duration trade-off, recovering the exact Pareto front where a ground truth exists. &
Ch.~\ref{chap:results} \\
\addlinespace

Role-conditioned action space &
Extending actions to event--role pairs exposes resource-allocation strategies at event-level granularity, unavailable to pool-level approaches. &
Ch.~\ref{chap:results}, \ref{chap:discussion} \\
\addlinespace

Component necessity (ablation) &
Isolates the contribution of shielding and learning, showing that the combination outperforms either component alone and that the observed optimisation quality is attributable to \ac{PPO} rather than to the restriction of the action space. &
Ch.~\ref{chap:internal_validation} \\
\addlinespace

External validity &
The compliance gap and the exact Pareto recovery both replicate on six mined \ac{DCR} graphs of independent provenance, never used to design the method; the same benchmark grid also exposes the scale at which this combination's advantage over a non-learning shielded baseline disappears. &
Ch.~\ref{chap:external_validation} \\

\bottomrule
\end{tabularx}
\end{table}

\paragraph{Main research question.} The main research question posed in Section~\ref{sec:research_questions} was:

\begin{quote}
\emph{How can \ac{RL} optimise declarative business-process execution while guaranteeing compliance with \ac{DCR} constraints?}
\end{quote}

Taken together, the components summarised in Table~\ref{tab:mrq_answers} constitute a direct answer: \ac{RL} can optimise declarative business-process execution while guaranteeing compliance by modelling \ac{DCR} execution as a shielded sequential decision problem, scalarising competing objectives into a unified reward signal, and extending the action space to joint event--role decisions. The framework does not replicate the exactness guarantees of symbolic solvers, and the learning component's advantage over structural restriction alone is bounded by the size of the action space it must search without guidance, but within that range it provides a scalable alternative that preserves formal compliance throughout training, generalises to process structures never seen during its development, and extends naturally to resource-aware optimisation scenarios.

% ------------------------------------------------------------
\section{Contributions}
\label{sec:contributions}

Section~\ref{sec:discussion_prior_work} showed that the cost--duration trade-off recovered in this thesis is not specific to reinforcement learning: it reappears, in the same qualitative direction, across three independent approaches that share no data, algorithm, or process formalism. What distinguishes the present work is not the trade-off itself, but three properties of how it is obtained:

\begin{description}

    \item[C1 --- Formal compliance guarantee.] Compliance is a hard structural property rather than a statistical outcome. Huang et al.\ operate on imperative models where the process graph itself prevents invalid transitions; Kirchdorfer et al.\ enforce compliance statistically through simulation fidelity. The present framework is the only one of the three that operates on a \emph{declarative} model, where the valid trace space is a non-trivial subset of all possible executions, and still guarantees that every generated trace is formally compliant with \ac{DCR} constraints throughout training, not only at convergence (Section~\ref{sec:discussion_compliance}).

    \item[C2 --- Event-level optimisation granularity.] The optimisation operates at the level of individual events rather than resource pools. This exposes the role-sensitivity structure discussed in Section~\ref{sec:discussion_roles}, identifying which specific activities drive the cost--duration trade-off and which are role-indifferent, an insight that pool-level approaches cannot produce.

    \item[C3 --- Native multi-objective optimisation.] Huang et al.\ require a separate training run per objective. The scalarised reward produces a full family of policies from a single weight sweep, allowing any point on the trade-off surface to be reached by adjusting $(\alpha, \beta)$ without redesigning the reward function or retraining from a different objective formulation.

\end{description}

These three contributions are evaluated internally against ablated variants of the same method (Chapter~\ref{chap:internal_validation}) and externally against six \ac{DCR} graphs never used to design it (Chapter~\ref{chap:external_validation}), establishing both that the components are individually necessary and that the resulting behaviour generalises beyond the two benchmarks used to develop the approach.

% ------------------------------------------------------------
\section{Future Work}
\label{sec:future_work}

The limitations identified in Section~\ref{sec:limitations} point to concrete directions for extending this work.

\paragraph{Action masking for large action spaces.} The single limitation with the most direct fix is the action-space ceiling identified in Chapter~\ref{chap:external_validation}: beyond approximately 42 event--role pairs, the unmasked policy's accept rate collapses because it keeps sampling actions that the shield blocks rather than only the currently enabled subset. Restricting the policy's sampling distribution directly, rather than only filtering its output, is a well-understood technique in the action-masking literature and was out of scope for this thesis only because of time, not because of any expected technical obstacle. Re-running the external validation grid with action masking enabled would test directly whether the learning component's advantage over the shielding-only baseline extends past the current ceiling.

\paragraph{Beyond linear scalarisation.} The reward design relies on linear scalarisation of cost and duration, which cannot recover Pareto points in non-convex regions of the objective space. Both benchmarks in this thesis happen to have convex fronts, so the limitation is not visible in the reported results, but it would become relevant for processes with more complex trade-off geometries. Hypervolume-indicator-based reward shaping, or explicitly multi-objective RL algorithms that maintain a population of policies rather than one policy per weight pair, are natural next steps, and would also reduce the need to retrain once per $(\alpha,\beta)$ configuration.

\paragraph{Algorithmic comparison.} All experiments use \ac{PPO} as the sole learning algorithm. Repeating the headline weight sweep with other policy-gradient or value-based methods would clarify whether the brevity bias, the convergence speed to full compliance, and the Pareto coverage reported here are properties of the proposed environment and reward design, or specific to \ac{PPO}.

\paragraph{Stochastic, continuous-time, and concurrent process models.} Of all the directions listed in this section, this is the one with the most direct bearing on industrial applicability. Duration is currently a per-event deterministic scalar, and cost models are synthetic, taken from benchmark definitions or fixed role multipliers; real processes instead exhibit resource contention, waiting times, concurrent activity execution, and stochastic, context-dependent costs, none of which the present framework represents. Integrating a stochastic or continuous-time simulation layer, similar in spirit to Doumeni's AgentSimulator-based approach (Section~\ref{sec:discussion_doumeni}), while retaining the \ac{DCR} shield for compliance, would not be a minor refinement: it would test whether the policy learned under deterministic assumptions degrades gracefully or breaks down once durations and costs vary at execution time, which is precisely the gap between the synthetic benchmarks used here and a real deployment.

\paragraph{Tighter Pareto verification on larger benchmarks.} The Loan Application Pareto front is approximated from sampled accepting episodes rather than verified against an exact reference, unlike the Expense Report benchmark. Extending the CSP+COP ground-truth methodology to the Loan Application's larger action space, even at higher computational cost, would establish whether the 26-point front reported here is close to exhaustive or whether low-probability Pareto-optimal traces are being missed.

\paragraph{Finer-grained robustness.} The $n=4$ seed study (Section~\ref{sec:policy_generalisation}) establishes that the headline Pareto front, illegal-action ratio, and acceptance rate are stable across seeds, but the role-specialisation categories of Section~\ref{sec:discussion_roles} were not re-evaluated per seed, and hyperparameters were fixed from preliminary runs rather than tuned per benchmark. Re-running the role-specialisation analysis across seeds, and a systematic hyperparameter search per benchmark, would tighten the confidence interval around both findings.

\paragraph{Generalisation beyond the benchmarks studied.} Chapter~\ref{chap:external_validation} already tests generalisation across six mined \ac{DCR} graphs, but all evaluation in this thesis remains within the DCR-js ecosystem and synthetic or mined cost/duration assignments. Applying the framework to a real industrial event log with measured (rather than assigned) resource costs, and to declarative formalisms other than \ac{DCR} where a comparable shielding mechanism could be defined, would be the strongest test of whether the contributions identified in Section~\ref{sec:contributions} generalise beyond this thesis's specific modelling choices.

% ------------------------------------------------------------
\section{Closing Remark}
\label{sec:closing_remark}

The central finding of this thesis is that compliance and optimisation do not have to compete for the same learning signal. By enforcing \ac{DCR} constraints structurally, through a shield, rather than statistically, through penalties the agent must learn to avoid, reinforcement learning is left free to focus entirely on the cost--duration trade-off, and it does so effectively: recovering exact Pareto fronts where they are known, well-populated approximations where they are not, and resource-allocation patterns that a pool-level method could not expose. The same architecture also makes its own limits legible, the action-space ceiling in Chapter~\ref{chap:external_validation} is not a failure of the method so much as a precise description of where its current implementation needs an additional mechanism, action masking, to keep working. That combination, a guarantee that holds unconditionally and a boundary that is known rather than assumed, is what this thesis offers as a foundation for compliance-aware process optimisation going forward.
